\section{Related Work}\label{sec:related}

Research on multi-agent systems has established several ways of distributing complex decision processes across interacting components. Classical agent-oriented work emphasizes autonomy, reactivity, pro-activeness, and social interaction as fundamental properties of agents ~\cite{wj95}. More recent LLM-based systems operationalize these ideas through structured collaboration. AutoGen, for example, uses conversational interaction as a mechanism for coordinating multiple agents ~\cite{autogen}, while MetaGPT organizes agents around specialized roles and standardized workflows ~\cite{metagpt}. These approaches show that dividing responsibility can structure complex tasks, but behavior depends also on how information and intermediate decisions move between agents.

This distinction becomes important when architectural organization is evaluated empirically. Studies comparing different agent configurations have reported task-dependent rather than universally superior arrangements ~\cite{kim2026artifact,kim2026capable}. In addition, evaluations centered on final task artifacts can obscure the internal mechanism responsible for an observed difference: a change in outcome may arise from role specialization, from the information exchanged between agents, or from the coordination procedure that determines how their outputs are combined. Recent analyses of multi-agent systems have similarly identified failures involving specification, inter-agent alignment, and verification ~\cite{mast}, indicating that the organization of interaction is itself an important part of the system being evaluated. The question is therefore not simply whether several agents can collectively solve a task, but whether the architectural separation of responsibilities changes the decisions produced when the surrounding information and decision conditions are held constant.

We treat adaptive interviewing as a sequential decision problem in which accumulated session evidence determines subsequent interactions. Reinforcement learning is one framework for mapping an observed state to actions over a sequence of decisions; our system uses PPO~\cite{ppo} as implemented in Stable-Baselines3~\cite{sb3}. In a technical interview, however, the information relevant to longitudinal assessment progression need not be identical to the information used for feedback, hints, or follow-up behavior. A centralized controller can therefore couple responsibilities that operate on different aspects of an evolving interaction. When the controller acts on a compact representation of a longer session, an individual decision also affects the state encountered at later turns, making the resulting trajectory an important object of analysis rather than merely a sequence of independent actions.

Separating assessment progression from coaching-related responsibilities may change the organization and ownership of state without changing the final decisions, whereas introducing additional decision sources can alter the trajectory through the mechanism used to admit and reconcile their proposals. The two effects are easy to conflate when decomposition and coordination are introduced simultaneously, so a controlled evaluation needs to distinguish them.

Safety-oriented reinforcement learning provides the complementary principle that adaptive decisions need not possess unrestricted authority. Shielding and supervisory approaches place external constraints around learned policies so that unsafe actions can be rejected, overridden, or restricted before they reach the environment~\cite{shield,safe}. This separation between policy generation and authoritative enforcement becomes particularly relevant when additional stateful components are introduced. A component may maintain information relevant to its role and recommend an action without being permitted to modify the authoritative state directly. Such a distinction allows proposal generation, validation, arbitration, and final commitment to be treated as separate responsibilities rather than allowing every participating component to act as a state owner.

Once several components can contribute recommendations, the system must determine which information each component owns, what authority accompanies a proposal, how proposals are reconciled with the underlying policy, and which constraints remain authoritative when they conflict. These questions require the internal decision path to be observable if their individual effects are to be separated experimentally.

The literature consequently provides the principal building blocks for distributed responsibility, sequential adaptation, and supervisory constraint enforcement, while leaving open an important empirical distinction between \textbf{coordination and specialization}. Distributed state and role separation may change the architecture without changing matched decisions, whereas coordination and arbitration may directly affect the trajectory through which decisions are produced. Authority boundaries add a further dimension by determining whether those distributed components remain subordinate to a single authoritative state owner. The present study uses this distinction as the basis for its architecture: specialized components maintain bounded private state and communicate constrained proposals, while authoritative commitment remains centralized. The architectural structure required to make those distinctions explicit is described next.
